\documentclass[11pt,a4paper]{article}
\usepackage[T1]{fontenc}
\usepackage{amsmath,amssymb,geometry}
\geometry{margin=28mm}
\title{EPPA--III: short two-point proof of the double-cover step}
\author{Independent validation note; manuscript unchanged}
\date{9 October 2026}
\begin{document}
\maketitle

\paragraph{Statement.}
Let \(G\) be a finite non-homogeneous graph and
\(H\) an EPPA witness for an induced copy of \(G\).
Suppose an \(\operatorname{Aut}(H)\)-invariant
system of imprimitivity partitions \(H\) into pairs
\(\{x_0,x_1\}\), one for each \(x\in G\), with
\(G=\{x_0:x\in V(G)\}\).
Then all edges between different pairs are
exactly those of the Taylor double of \(G\);
edges within the pairs are either all present or
all absent. Furthermore, \(\operatorname{Aut}(H)\)
induces a 2-transitive group on the pairs and
acts transitively on ordered inter-pair edges
and on ordered inter-pair nonedges.

\paragraph{Proof of the structure.}
Because the pair system is invariant, every
automorphism of \(H\) induces a permutation
of the pairs, and if it maps \(x_0\) to \(y_0\)
it must map \(x_1\) to \(y_1\).
Every ordered edge of \(G\) can be carried by
an EPPA extension to every other ordered edge,
and likewise for ordered nonedges.
Consequently, the adjacency matrix between
two different pairs depends only on whether
\(x\sim_G y\). Writing its rows and columns in
the order \(0,1\), the matrices have the form
\[
  A_E=\begin{pmatrix}1&c_E\\c_E&e_T\end{pmatrix},
  \qquad
  A_N=\begin{pmatrix}0&c_N\\c_N&n_T\end{pmatrix},
  \qquad c_E,c_N,e_T,n_T\in\{0,1\}.
\]
The off-diagonal entries agree because EPPA
extends the transposition of the ordered pair.
Since \(G\) is non-homogeneous, it has both
edges and nonedges.

Let \(m_E=1+2c_E+e_T\) and \(m_N=2c_N+n_T\)
be the numbers of edges between the two pairs.
If \(m_E\ne m_N\), every automorphism of \(H\)
preserves the edge/nonedge distinction on the
underlying set of pairs. Thus every partial
automorphism of \(G\), extended to \(H\) and
then projected to the invariant pair system,
extends to an automorphism of \(G\).
That would make \(G\) homogeneous. Hence
\(m_E=m_N\).

The four binary parameters now have precisely
three possible configurations:
\[
\begin{array}{c|cc}
  &A_E&A_N\\ \hline
  \text{one edge}&
  \bigl(\begin{smallmatrix}1&0\\0&0\end{smallmatrix}\bigr)&
  \bigl(\begin{smallmatrix}0&0\\0&1\end{smallmatrix}\bigr)\\[2pt]
  \text{two edges}&
  \bigl(\begin{smallmatrix}1&0\\0&1\end{smallmatrix}\bigr)&
  \bigl(\begin{smallmatrix}0&1\\1&0\end{smallmatrix}\bigr)\\[2pt]
  \text{three edges}&
  \bigl(\begin{smallmatrix}1&1\\1&0\end{smallmatrix}\bigr)&
  \bigl(\begin{smallmatrix}0&1\\1&1\end{smallmatrix}\bigr)
\end{array}
\]
In the first configuration, the \emph{unique
edge} between every two different pairs joins
vertices in the same layer. In the third,
the \emph{unique nonedge} between every two
different pairs likewise joins vertices in
the same layer. An automorphism that flips
one pair but not another would turn that
unique edge (respectively nonedge) into a
cross-layer edge (respectively nonedge),
which is impossible. Therefore every
automorphism flips either \emph{all} pairs
or \emph{none}. Any EPPA extension of a
nonempty partial automorphism \(G\to G\)
must fix the bottom layer setwise, and
hence restricts to an automorphism of \(G\).
This again contradicts non-homogeneity.

Only the middle configuration remains,
which is exactly the Taylor double between
different fibres. The mate adjacency
\(x_0\sim x_1\) is constant by the one-point
EPPA extensions, so it can be uniformly
added or omitted.

\paragraph{The orbit conclusion.}
Let \(K\) be the group induced by
\(\operatorname{Aut}(H)\) on the pairs.
Two-point EPPA shows that \(K\) is
transitive separately on ordered edges and
ordered nonedges of \(G\). If these were
distinct \(K\)-orbits, then \(K\le
\operatorname{Aut}(G)\), forcing \(G\)
to be homogeneous as above. Hence the
two orbits coalesce: \(K\) is 2-transitive.
Every pair of distinct fibres in the
Taylor-double pattern has two edges
and two nonedges between them. The
deck involution exchanging \(x_0,x_1\)
simultaneously for every \(x\) exchanges
the two edges and the two nonedges.
Combining it with 2-transitivity of \(K\)
proves transitivity on ordered edges and
ordered nonedges connecting different fibres.

\paragraph{Consequences.}
This replaces the maximal-partial-automorphism
argument in Claim 3.5 and the ensuing counting
in Claim 3.6 with an elementary two-point
orbit calculation. It also gives the Case~2
double-cover identification in Lemma 3.12
\emph{without using regularity or unequal
clique sizes}, once the pair system is known.
For Lemma 3.3, Claim 3.4 remains the separate
step establishing that invariant pair system.

\paragraph{Verification boundary.}
The three binary configurations are exhaustively
checked and encoded in Lean in
\texttt{EPPAIII/DoubleCover/PairPatterns.lean}.
The logical implication from an invariant pair
system and EPPA to the uniform matrices, and
the layer-rigidity implication, remain to be
fully formalized. The source manuscript has
not been changed.
\end{document}
